\documentclass[14pt,a4paper,oneside]{extbook}
\usepackage[utf8]{vietnam}
\usepackage{amsmath,amssymb,mathrsfs,fancyhdr,enumerate,multirow,makecell}
\usepackage{tkz-euclide,tikz-3dplot,tikz,tkz-tab}
\usetikzlibrary{arrows.meta,calc,intersections,angles,quotes,backgrounds,shapes.geometric,patterns,through,shadows,positioning}
\usepackage{yhmath}
\usepackage[top=0.7cm, bottom=0.7cm, left=1.0cm, right=1.0cm, headheight=16pt, headsep=3mm, footskip=5mm]{geometry}
\usepackage[hidelinks,unicode]{hyperref}
\usepackage{qrcode}
\usepackage[many]{tcolorbox}
\usepackage{fontawesome}

\pagestyle{fancy}
\fancyhf{}
\fancyhead[L]{\footnotesize\bfseries\color{teal!85!black}\faBook~SỔ TAY KIẾN THỨC TOÁN 5}
\fancyhead[R]{\footnotesize\bfseries\color{blue!85!black}Lớp Toán Thầy Lĩnh -- loptoanthaylinh.vn}
\fancyfoot[C]{\footnotesize\bfseries\color{gray!80!black}-- Sổ tay kiến thức: Bài 68. Ôn tập số tự nhiên, phân số, số thập phân --}
\renewcommand{\headrulewidth}{0.4pt}

\begin{document}

% --- BANNER TIÊU ĐỀ BÀI HỌC ---
\noindent
\begin{tcolorbox}[
  enhanced,
  colback=teal!8!white,
  colframe=teal!80!black,
  arc=3mm,
  boxrule=1.1pt,
  width=\linewidth,
  halign=center,
  top=2.2mm, bottom=2.2mm,
  left=3mm, right=3mm,
  drop fuzzy shadow
]
  {\large\bfseries\color{teal!95!black}\faBookmark\ BÀI 68. ÔN TẬP SỐ TỰ NHIÊN, PHÂN SỐ, SỐ THẬP PHÂN}\\[2pt]
  {\scriptsize\color{black!75}\itshape Tổng ôn cấu tạo, dấu hiệu chia hết, quy đồng phân số và so sánh số thập phân trên 1 trang A4}
\end{tcolorbox}

\vspace{3.2mm}
% --- KHỐI 1: SỐ TỰ NHIÊN VÀ PHÂN SỐ ---
\noindent
\begin{minipage}[t]{0.52\textwidth}
  \begin{tcolorbox}[
    enhanced,
    title={\footnotesize\bfseries\color{white}\faHashtag\ 1. Số tự nhiên \& Dấu hiệu chia hết},
    colback=blue!2!white,
    colframe=blue!75!black,
    arc=2.5mm, boxrule=0.8pt, top=1.8mm, bottom=1.8mm, left=2.2mm, right=2.2mm
  ]
    \footnotesize
    $\bullet$ \textbf{Hàng và lớp:} Mỗi lớp gồm $3$ hàng (Đơn vị, Chục, Trăm).\\
    $\bullet$ \textbf{Chia hết cho 2, 5:} Tận cùng là chẵn ($0, 2, 4, 6, 8$) hoặc $0, 5$.\\
    $\bullet$ \textbf{Chia hết cho 3, 9:} Tổng các chữ số chia hết cho $3, 9$.
  \end{tcolorbox}
\end{minipage}%
\hfill
\begin{minipage}[t]{0.46\textwidth}
  \begin{tcolorbox}[
    enhanced,
    title={\footnotesize\bfseries\color{white}\faSliders\ 2. Phân số \& Số thập phân},
    colback=teal!2!white,
    colframe=teal!75!black,
    arc=2.5mm, boxrule=0.8pt, top=1.8mm, bottom=1.8mm, left=2mm, right=2mm
  ]
    \footnotesize
    $\bullet$ \textbf{Phân số:} Rút gọn về tối giản, quy đồng mẫu số chung.\\
    $\bullet$ \textbf{Số thập phân:} So sánh phần nguyên trước; nếu bằng nhau thì so sánh lần lượt phần mười, phần trăm, phần nghìn...
  \end{tcolorbox}
\end{minipage}

\vspace{3mm}

% --- KHỐI 3: SƠ ĐỒ TƯ DUY ---
\noindent
\begin{tcolorbox}[
  enhanced,
  title={\footnotesize\bfseries\color{white}\faSitemap\ Sơ đồ tư duy: Bản đồ 3 hệ thống số trong Toán tiểu học},
  colback=yellow!2!white,
  colframe=teal!85!black,
  arc=3mm, boxrule=0.9pt, top=1.8mm, bottom=1.8mm, left=1.5mm, right=1.5mm,
  halign=center
]
  \begin{tikzpicture}[
    scale=0.92,
    every node/.style={transform shape},
    root/.style={rectangle, rounded corners=5pt, fill=teal!85!black, text=white, font=\bfseries\footnotesize, inner sep=5pt, align=center},
    b1/.style={rectangle, rounded corners=4pt, fill=blue!12, draw=blue!75!black, text=blue!90!black, font=\bfseries\scriptsize, inner sep=4.5pt, align=center},
    b2/.style={rectangle, rounded corners=4pt, fill=purple!12, draw=purple!75!black, text=purple!90!black, font=\bfseries\scriptsize, inner sep=4.5pt, align=center},
    b3/.style={rectangle, rounded corners=4pt, fill=green!15, draw=green!65!black, text=green!80!black, font=\bfseries\scriptsize, inner sep=4.5pt, align=center}
  ]
    \node[root] (core) at (0, 0.85) {\faBook\ CÁC TẬP HỢP SỐ};
    \node[b1] (n1) at (-5.0, -0.45) {\textbf{SỐ TỰ NHIÊN}\\[2pt]Hàng và Lớp, đọc viết,\\Dấu hiệu chia hết ($2, 3, 5, 9$)};
    \node[b2] (n2) at (0, -0.65) {\textbf{PHÂN SỐ \& HỖN SỐ}\\[2pt]Rút gọn PSTG, quy đồng,\\So sánh (phần bù, phần hơn)};
    \node[b3] (n3) at (5.0, -0.45) {\textbf{SỐ THẬP PHÂN}\\[2pt]Nguyên , Thập phân,\\Số thập phân bằng nhau, làm tròn};

    \draw[teal!85!black, thick, -{Stealth[length=2mm]}] (core) -| (n1);
    \draw[teal!85!black, thick, -{Stealth[length=2mm]}] (core) -- (n2);
    \draw[teal!85!black, thick, -{Stealth[length=2mm]}] (core) -| (n3);
  \end{tikzpicture}
\end{tcolorbox}

\vspace{3mm}

% --- KHỐI 4: CHUYỂN ĐỔI & BẪY TRẮC NGHIỆM ---
\noindent
\begin{minipage}[t]{0.54\textwidth}
  \begin{tcolorbox}[
    enhanced,
    title={\footnotesize\bfseries\color{white}\faExchange\ Quy tắc chuyển đổi ba dạng số},
    colback=green!3!white,
    colframe=green!60!black,
    arc=2.5mm, boxrule=0.8pt, top=1.8mm, bottom=1.8mm, left=2.2mm, right=2.2mm
  ]
    \footnotesize
    $\bullet$ Phân số ra STP: Lấy tử chia cho mẫu ($\frac{3}{4} = 3 : 4 = 0,75$).\\
    $\bullet$ Hỗn số ra phân số: $\text{Tử} = \text{Nguyên} \times \text{Mẫu} + \text{Tử cũ}$ ($2\frac{1}{3} = \frac{7}{3}$).\\
    $\bullet$ Số tự nhiên viết dưới dạng phân số có mẫu là $1$ ($a = \frac{a}{1}$).
  \end{tcolorbox}
\end{minipage}%
\hfill
\begin{minipage}[t]{0.44\textwidth}
  \begin{tcolorbox}[
    enhanced,
    title={\footnotesize\bfseries\color{white}\faExclamationTriangle\ Mẹo tránh bẫy trắc nghiệm},
    colback=red!2!white,
    colframe=red!70!black,
    arc=2.5mm, boxrule=0.8pt, top=1.8mm, bottom=1.8mm, left=2.2mm, right=2.2mm
  ]
    \footnotesize
    $\bullet$ Bỏ số $0$ ở tận cùng phần thập phân giá trị không đổi ($0,50 = 0,5$), nhưng số $0$ ở giữa \textbf{tuyệt đối không được bỏ} ($4,05 \neq 4,5$).\\
    $\bullet$ Nhiều chữ số TP hơn chưa chắc lớn hơn ($0,125 < 0,5$).
  \end{tcolorbox}
\end{minipage}

\end{document}
